\documentclass[12pt,a4paper]{article}
\usepackage[utf8]{inputenc}
\usepackage[T1]{fontenc}
\usepackage[spanish]{babel}
\usepackage{amsmath}
\usepackage{amsfonts}
\usepackage{amssymb}
\usepackage{graphicx}
\usepackage{geometry}
\usepackage{setspace}
\usepackage{titlesec}
\usepackage{fancyhdr}
\usepackage{cite}
\usepackage{url}
\usepackage{hyperref}
\usepackage{booktabs}

% Configuración de página
\geometry{left=3cm,right=2.5cm,top=2.5cm,bottom=2.5cm}
\onehalfspacing

% Configuración de títulos
\titleformat{\section}{\normalfont\Large\bfseries}{\thesection}{1em}{}
\titleformat{\subsection}{\normalfont\large\bfseries}{\thesubsection}{1em}{}
\titleformat{\subsubsection}{\normalfont\normalsize\bfseries}{\thesubsubsection}{1em}{}

% Información del documento
\title{\textbf{SISTEMA AVANZADO DE RECOMENDACIONES MUSICALES BASADO EN ANÁLISIS SEMÁNTICO DE LETRAS Y CLUSTERING OPTIMIZADO}\\
\vspace{0.5cm}
\large{Anteproyecto de Tesis de Grado}\\
\large{Ingeniería Informática}}

\author{
\textbf{Estudiante 1:} San Miguel, Santiago \\
\textbf{Estudiante 2:} Cordoba, Thiago \\
\\
\textbf{Tutor:} Rivera Bernsdorff, Fernando Lucas \\
\\
\textbf{Universidad:} Universidad Católica de Salta \\
\textbf{Facultad:} Facultad de Ingeniería \\
\textbf{Carrera:} Ingeniería Informática
}

\date{Septiembre 2025}

\begin{document}

\maketitle
\thispagestyle{empty}

\newpage

% =====================================================
% RESUMEN DEL PROYECTO
% =====================================================

\section*{Resumen del Proyecto}

Este proyecto aborda una problemática interesante y poco explorada en los sistemas de recomendación musical: la integración del análisis semántico de letras musicales como dimensión complementaria a las características acústicas tradicionales. La propuesta revoluciona el enfoque convencional mediante la implementación de vectorización contextual avanzada con modelos BERT, transformando el contenido lírico en representaciones vectoriales de 384 dimensiones que capturan la esencia temática, emocional y narrativa de las canciones. El sistema desarrolla técnicas innovadoras de clustering optimizado que purifican y estructuran tanto el espacio musical de características acústicas Spotify como el espacio semántico de embeddings textuales, creando una arquitectura híbrida única que comprende verdaderamente qué significa una canción más allá de cómo suena. A través del procesamiento inteligente de un dataset curado de 18,000 canciones con metadatos completos y letras verificadas, el proyecto demuestra que las palabras importan tanto como la música, estableciendo un nuevo paradigma en Music Information Retrieval donde las recomendaciones trascienden la similitud acústica para abrazar la coherencia temática y la relevancia semántica profunda.

\newpage
\tableofcontents
\newpage

% =====================================================
% 1. INTRODUCCIÓN
% =====================================================
\section{Introducción}

Los sistemas de recomendación musical representan uno de los pilares tecnológicos más críticos en la industria del entretenimiento digital contemporáneo, procesando millones de consultas diarias en plataformas como Spotify, Apple Music y YouTube Music. Sin embargo, los enfoques actuales presentan una limitación fundamental: se basan exclusivamente en características acústicas extraídas automáticamente de las grabaciones musicales (tempo, energía, valencia emocional, instrumentalidad) o en patrones de comportamiento de usuarios similares, ignorando completamente el contenido semántico de las letras musicales.

Esta limitación resulta en recomendaciones que, aunque musicalmente coherentes desde una perspectiva técnica, pueden ser temáticamente incongruentes o contextualmente inapropiadas para las preferencias específicas del usuario. Por ejemplo, un sistema tradicional podría recomendar canciones con características acústicas similares pero con contenido lírico completamente opuesto (canciones alegres musicalmente pero con letras melancólicas), degradando significativamente la experiencia de usuario y limitando la precisión de las recomendaciones.

Las letras musicales constituyen una dimensión informacional rica y complementaria que contiene información semántica, temática, emocional y contextual que no puede inferirse directamente desde características acústicas. La integración sistemática de análisis de contenido lírico representa una oportunidad técnica y científica de considerable relevancia que ha permanecido relativamente inexplorada en la literatura académica especializada.

Los avances recientes en procesamiento de lenguaje natural, particularmente los modelos transformer como BERT (Bidirectional Encoder Representations from Transformers), proporcionan capacidades sin precedentes para generar representaciones vectoriales contextuales que capturan relaciones semánticas complejas, dependencias contextuales y patrones temáticos en textos musicales. Esta tecnología, combinada con técnicas optimizadas de clustering musical, ofrece el potencial de desarrollar sistemas de recomendación substancialmente superiores en términos de precisión, relevancia temática y satisfacción de usuario.

\subsection{Justificación de la Propuesta}

La justificación de este proyecto se fundamenta en tres pilares principales que evidencian su relevancia técnica, científica y práctica:

\textbf{Brecha Tecnológica Identificada}: La revisión sistemática de la literatura científica y el análisis de sistemas comerciales actuales revela una ausencia notable de sistemas que integren efectivamente análisis semántico avanzado de letras con clustering musical optimizado. Los pocos trabajos existentes utilizan técnicas rudimentarias de análisis textual que resultan inadecuadas para capturar la complejidad semántica del lenguaje lírico contemporáneo.

\textbf{Viabilidad Técnica Demostrada}: Los avances en modelos de lenguaje contextuales y la disponibilidad de datasets musicales con metadatos completos proporcionan la infraestructura técnica necesaria para abordar este desafío. La propuesta aprovecha tecnologías maduras (BERT, clustering optimizado) aplicándolas de manera innovadora al dominio musical específico.

\textbf{Impacto Potencial Significativo}: El desarrollo exitoso de un sistema híbrido que combine efectivamente información musical y semántica posee implicaciones directas para la industria del streaming musical, con potencial de mejoras substanciales en engagement de usuarios, precisión de recomendaciones, y desarrollo de aplicaciones especializadas en contextos terapéuticos, educativos o de entretenimiento específico.

% =====================================================
% 2. OBJETIVOS
% =====================================================
\section{Objetivos}

\subsection{Objetivo General}

Desarrollar e implementar un sistema avanzado de recomendaciones musicales que integre análisis semántico de letras mediante embeddings BERT con clustering optimizado de características acústicas, demostrando mejoras cuantificables en precisión y relevancia de recomendaciones comparado con enfoques unimodales tradicionales.

\subsection{Objetivos Específicos}

\begin{enumerate}
    \item \textbf{Implementar vectorización semántica de letras musicales}: Desarrollar un sistema de procesamiento que utilice modelos BERT pre-entrenados para generar representaciones vectoriales contextuales de letras musicales, optimizado para datasets de gran escala.
    
    \item \textbf{Optimizar clustering musical mediante purificación híbrida}: Diseñar e implementar técnicas de purificación que mejoren significativamente las métricas de calidad de clustering musical a través de eliminación inteligente de outliers y selección de características discriminativas.
    
    \item \textbf{Desarrollar arquitectura de fusión multimodal}: Crear un sistema que combine efectivamente información de espacios vectoriales musical (12D) y semántico (384D) mediante técnicas de normalización y ponderación científicamente validadas.
    
    \item \textbf{Validar experimentalmente el sistema híbrido}: Demostrar mediante experimentación controlada que el enfoque multimodal produce mejoras estadísticamente significativas en métricas de recomendación versus sistemas baseline unimodales.
    
\end{enumerate}

% =====================================================
% 3. ALCANCE DEL PROYECTO
% =====================================================
\section{Alcance del Proyecto}

Este proyecto ambicioso abarca el desarrollo completo de un ecosistema tecnológico innovador que revoluciona la forma de entender y recomendar música. Desde la transformación de datos masivos hasta la creación de interfaces intuitivas, cada componente está diseñado para demostrar que las letras musicales contienen información tan valiosa como las características acústicas tradicionales:

\begin{itemize}
    \item \textbf{Análisis exhaustivo y purificación del dataset original}: Procesamiento inteligente del dataset inicial de más de 1.2 millones de canciones, aplicando técnicas avanzadas de limpieza, normalización de metadatos, eliminación de duplicados, y verificación de integridad referencial para obtener un conjunto curado de 18,000 canciones de alta calidad con letras completas verificadas mediante múltiples APIs especializadas.
    
    \item \textbf{Sistema revolucionario de vectorización BERT para letras musicales}: Implementación completa de procesamiento de lenguaje natural que transforma el contenido lírico en representaciones vectoriales contextuales de 384 dimensiones, capturando significados profundos, relaciones semánticas, contextos emocionales, y narrativas temáticas que los enfoques tradicionales no pueden detectar.
    
    \item \textbf{Algoritmos innovadores de clustering optimizado dual}: Desarrollo de técnicas de purificación híbrida que estructuran tanto el espacio musical de características acústicas (12D) como el espacio semántico de embeddings textuales (384D), creando agrupamientos interpretativos y cohesivos que revelan patrones ocultos en ambas modalidades.
    
    \item \textbf{Motor de fusión multimodal científicamente validado}: Arquitectura híbrida que combina inteligentemente información musical y semántica mediante técnicas de normalización, ponderación adaptiva, y fusión tardía, generando recomendaciones que comprenden tanto el sonido como el significado de las canciones.
    
    \item \textbf{Sistema de evaluación experimental comprensivo}: Suite completa de métricas especializadas para validación de sistemas multimodales, incluyendo análisis de precisión, diversidad, coherencia cross-modal, y benchmarking comparativo contra enfoques unimodales tradicionales.
    
    \item \textbf{Interfaz de usuario amigable e intuitiva}: Desarrollo potencial de una aplicación web interactiva que permita a usuarios explorar recomendaciones, visualizar clusters temáticos, y comprender las razones detrás de cada sugerencia musical mediante explicaciones automáticas basadas en contenido lírico y características acústicas.
    
    \item \textbf{Documentación técnica y científica completa}: Especificaciones detalladas de arquitectura, protocolos experimentales reproducibles, análisis de resultados, y contribuciones metodológicas que permitan extensión y replicación del trabajo por parte de la comunidad académica.
\end{itemize}


% =====================================================
% 4. MARCO TEÓRICO
% =====================================================
\section{Marco Teórico}

\subsection{Sistemas de Recomendación Musical}

Los sistemas de recomendación musical se clasifican en tres paradigmas principales: filtrado colaborativo (análisis de patrones de usuarios similares), filtrado basado en contenido (características intrínsecas de canciones), y sistemas híbridos que combinan múltiples enfoques. El presente proyecto se enfoca en el paradigma basado en contenido, específicamente en la extensión del análisis tradicional de características acústicas hacia la incorporación de información semántica textual.

\subsection{Clustering Musical y Métricas de Calidad}

El clustering musical busca identificar agrupamientos naturales en espacios de características que correspondan a conceptos musicales interpretables. Las métricas principales para evaluar calidad de clustering incluyen el Silhouette Score (coherencia intra-cluster vs separación inter-cluster), Calinski-Harabasz Index (ratio de dispersión), y Hopkins Statistic (evaluación de clustering tendency del dataset).

\subsection{Modelos BERT y Análisis Semántico}

BERT (Bidirectional Encoder Representations from Transformers) representa el estado del arte en modelos de lenguaje contextuales, utilizando arquitecturas de atención bidireccional para generar representaciones vectoriales que capturan relaciones semánticas complejas. Para aplicaciones musicales, BERT puede procesar letras completas y generar embeddings de dimensionalidad fija (384D para BERT-base) que preservan información semántica y temática.

\subsection{Fusión Multimodal}

La integración efectiva de información musical y textual requiere técnicas especializadas que aborden la asimetría dimensional (12D vs 384D), diferencias en distribución de datos, y optimización de pesos de combinación. Los enfoques de fusión tardía mediante combinación ponderada de scores de similaridad han demostrado efectividad en dominios similares.

% =====================================================
% 5. METODOLOGÍA DE TRABAJO Y DIVISIÓN DE RESPONSABILIDADES
% =====================================================
\section{Metodología de Trabajo y División de Responsabilidades}

\subsection{Organización del Equipo}

El proyecto se estructura como una colaboración técnica entre dos estudiantes de ingeniería informática, implementando especialización modular con integración continua. La división aprovecha la separación natural entre dominios técnicos (musical vs semántico) con fases de colaboración intensiva en análisis de datos, integración multimodal, y validación experimental.

\subsection{División Específica de Responsabilidades}

\textbf{Desarrollador A - Dominio Musical}: Responsable del pipeline completo de características musicales, implementación de algoritmos de clustering optimizado con purificación híbrida, desarrollo de métricas de evaluación de clustering, y motor de recomendación basado exclusivamente en características acústicas.

\textbf{Desarrollador B - Dominio Semántico}: Responsable del sistema de extracción de letras mediante múltiples APIs, implementación de vectorización BERT con optimizaciones de memoria, desarrollo de clustering semántico en alta dimensionalidad, y motor de recomendación basado en similaridad de embeddings.

\subsection{Fases de Colaboración Crítica}

\textbf{Análisis y Curación de Datos}: Definición conjunta de criterios de calidad, análisis exploratorio de distribuciones en ambos dominios, diseño de pipeline unificado de limpieza, y validación cruzada de integridad de datos.

\textbf{Integración Multimodal}: Diseño colaborativo de arquitectura de fusión, experimentación conjunta para optimización de pesos, validación de coherencia en recomendaciones híbridas, y desarrollo de suite de testing integral.

\textbf{Validación y Documentación}: Diseño conjunto de protocolos experimentales, análisis colaborativo de resultados, resolución de problemas de integración, y preparación de documentación técnica comprensiva.

% =====================================================
% 6. CRONOGRAMA
% =====================================================
\section{Cronograma}

Este proyecto se despliega estratégicamente en 6 fases interconectadas a lo largo de 8 meses intensivos de desarrollo, donde cada etapa construye sobre los logros de la anterior para crear un sistema verdaderamente revolucionario:

\textbf{FASE 1 - Análisis de Datos (6 semanas):}
\begin{itemize}
    \item Extracción y curación inteligente del dataset masivo inicial
    \item Transformación de 1.2M canciones brutas en 18K canciones verificadas
    \item Análisis exploratorio que revele los secretos ocultos en patrones musicales y semánticos
    \item Validación científica de la viabilidad del clustering en ambos dominios
\end{itemize}

\textbf{FASE 2 - Construcción de los Motores (8 semanas):}
\begin{itemize}
    \item Desarrollo paralelo e intensivo de los módulos musical y semántico
    \item Implementación del corazón BERT que entiende el significado de las letras
    \item Forjado de algoritmos de clustering que purifican y revelan estructura oculta
    \item Primeros motores de recomendación unimodales
\end{itemize}

\textbf{FASE 3 - Integración de Motores (4 semanas):}
\begin{itemize}
    \item Diseño de la arquitectura híbrida que unifica música y significado
    \item Implementación de módulo central que piensa en ambos lenguajes
    \item Testing exhaustivo del sistema
    \item Optimización de la comunicación entre mundos musical y semántico
\end{itemize}

\textbf{FASE 4 - Validación Científica (6 semanas):}
\begin{itemize}
    \item Experimentación rigurosa que demuestre la superioridad del enfoque híbrido
    \item Estudios que revelen la contribución de cada componente
    \item Refinamiento basado en evidencia empírica
\end{itemize}

\textbf{FASE 5 - Documentación (4 semanas):}
\begin{itemize}
    \item Documentación exhaustiva de la arquitectura planteada
    \item Análisis profundo de las contribuciones científicas logradas
\end{itemize}

% =====================================================
% 7. REFERENCIAS BIBLIOGRÁFICAS
% =====================================================
\section{Referencias Bibliográficas}

\begin{thebibliography}{99}

\bibitem{adomavicius2005toward}
Adomavicius, G., \& Tuzhilin, A. (2005). Toward the next generation of recommender systems: A survey of the state-of-the-art and possible extensions. \textit{IEEE transactions on knowledge and data engineering}, 17(6), 734-749.

\bibitem{devlin2018bert}
Devlin, J., Chang, M. W., Lee, K., \& Toutanova, K. (2018). Bert: Pre-training of deep bidirectional transformers for language understanding. \textit{arXiv preprint arXiv:1810.04805}.

\bibitem{schedule2018spotify}
Schedl, M., Zamani, H., Chen, C. W., Deldjoo, Y., \& Elahi, M. (2018). Current challenges and visions in music recommender systems research. \textit{International journal of multimedia information retrieval}, 7(2), 95-116.

\bibitem{yang2012machine}
Yang, Y. H., \& Chen, H. H. (2012). Machine recognition of music emotion: A review. \textit{ACM Transactions on Intelligent Systems and Technology (TIST)}, 3(3), 1-30.

\bibitem{mcfee2012million}
McFee, B., Bertin-Mahieux, T., Ellis, D. P., \& Lanckriet, G. R. (2012). The million song dataset challenge. \textit{Proceedings of the 21st international conference on World Wide Web}, 909-916.

\bibitem{hu2008collaborative}
Hu, Y., Koren, Y., \& Volinsky, C. (2008). Collaborative filtering for implicit feedback datasets. \textit{2008 Eighth IEEE International Conference on Data Mining}, 263-272.

\bibitem{celma2010music}
Celma, Ò. (2010). \textit{Music recommendation and discovery: the long tail, long fail, and long play in the digital music space}. Springer Science \& Business Media.

\bibitem{turnbull2008semantic}
Turnbull, D., Barrington, L., Torres, D., \& Lanckriet, G. (2008). Semantic annotation and retrieval of music and sound effects. \textit{IEEE Transactions on Audio, Speech, and Language Processing}, 16(2), 467-476.

\bibitem{zhang2019neural}
Zhang, S., Yao, L., Sun, A., \& Tay, Y. (2019). Deep learning based recommender system: A survey and new perspectives. \textit{ACM Computing Surveys}, 52(1), 1-38.

\bibitem{mesnage2011using}
Mesnage, C., Rafii, Z., \& Peeters, G. (2011). Using language models to improve audio music similarity. \textit{12th International Society for Music Information Retrieval Conference}, 475-480.

\end{thebibliography}

\end{document}